\subsection*{Параметры исследования}
\begin{table}[H]
    \centering
    \caption{Исходные данные для выполнения работы}
    \label{tab:initial_params}
    \begin{tblr}{
        colspec = { X[c,m] X[c,m] X[c,m] },
        hlines, vlines,
    }
        Тип звена & Номинальный коэф. передачи $k_H$ & Допустимая погрешность установления, \% \\
        ФНЧ 2-го порядка & 1 & 5 \\
    \end{tblr}
\end{table}

\subsection*{1. Исследование переходного процесса (Ступенчатое воздействие)}

\begin{table}[H]
    \centering
    \caption{Параметры звена для п. 1}
    \label{tab:params_step}
    \begin{tblr}{
        colspec = { Q[r,m] Q[l,m] Q[r,m] Q[l,m] },
    }
        $f_0 =$ & \_\_\_\_\_ Гц & $\beta =$ & \_\_\_\_\_ \\
    \end{tblr}
\end{table}

\begin{longtblr}[
    caption = {Результаты измерений при ступенчатом воздействии},
    label = {tab:step_response},
]{
    width = \linewidth,
    colspec = { Q[c,m] X[c,m] X[c,m] X[c,m] },
    hlines, vlines,
    rowhead = 1,
}
    № точки & Время $t_i$, мс & Вход $u_{\text{вх} i}$, В & Выход $u_{\text{вых} i}$, В \\
    1 & & & \\
    2 & & & \\
    3 & & & \\
    4 & & & \\
    5 & & & \\
    6 & & & \\
    7 & & & \\
    8 & & & \\
    9 & & & \\
    10 & & & \\
\end{longtblr}

\subsection*{2. Зависимость времени установления}

\begin{table}[H]
    \centering
    \caption{Зависимость $t_y = F(f_{0i})$ при $\beta = \text{const}$}
    \label{tab:settling_freq}
    \begin{tblr}{
        colspec = { Q[c,m] X[c,m] X[c,m] },
        hlines, vlines,
    }
        № опыта & Частота $f_{0i}$ (условн. ед. / Гц) & Время установления $t_y$, мс \\
        1 & & \\
        2 & & \\
        3 & & \\
        4 & & \\
    \end{tblr}
\end{table}

\begin{table}[H]
    \centering
    \caption{Зависимость $t_y = F(\beta_i)$ при $f_0 = \text{const}$}
    \label{tab:settling_beta}
    \begin{tblr}{
        colspec = { Q[c,m] X[c,m] X[c,m] },
        hlines, vlines,
    }
        № опыта & Коэф. демпфирования $\beta_i$ & Время установления $t_y$, мс \\
        1 & & \\
        2 & & \\
        3 & & \\
        4 & & \\
    \end{tblr}
\end{table}

\subsection*{3. Исследование синусоидального воздействия}

\begin{table}[H]
    \centering
    \caption{Параметры звена для п. 3}
    \label{tab:params_sine}
    \begin{tblr}{
        colspec = { Q[r,m] Q[l,m] Q[r,m] Q[l,m] },
    }
        $f_0 =$ & \_\_\_\_\_ Гц & $\beta =$ & \_\_\_\_\_ \\
    \end{tblr}
\end{table}

\begin{longtblr}[
    caption = {Результаты измерений при синусоидальном воздействии},
    label = {tab:sine_response},
]{
    width = \linewidth,
    colspec = { Q[c,m] X[c,m] X[c,m] X[c,m] },
    hlines, vlines,
    rowhead = 1,
}
    № точки & Время $t_i$, мс & Вход $u_{\text{вх} i}$, В & Выход $u_{\text{вых} i}$, В \\
    1 & & & \\
    2 & & & \\
    3 & & & \\
    4 & & & \\
    5 & & & \\
    6 & & & \\
    7 & & & \\
    8 & & & \\
    9 & & & \\
    10 & & & \\
\end{longtblr}